\chapter{Unit 6: Human Communities, Environmental Ethics, Laws and Movements}

\section{Section A: Multiple Choice Practice Questions (60 Curated Problems)}

\noindent\textbf{Q1. What does human population growth mean?} \hfill [Tag: Human population growth | Easy]
\begin{enumerate}[label=(\alph*),topsep=2pt,itemsep=1pt,partopsep=0pt,parsep=0pt]
    \item {Change in the age of one person}
    \item {Increase in the average height of people}
    \item {Increase in the number of people over time}
    \item {Movement of people between two cities}
\end{enumerate}
\vspace{2mm}

\noindent\textbf{Q2. Which factor directly contributes to rapid population growth?} \hfill [Tag: Human population growth | Easy]
\begin{enumerate}[label=(\alph*),topsep=2pt,itemsep=1pt,partopsep=0pt,parsep=0pt]
    \item {Lower birth rate}
    \item {Higher birth rate}
    \item {Higher death rate}
    \item {Greater emigration}
\end{enumerate}
\vspace{2mm}

\noindent\textbf{Q3. Which environmental problem commonly results from rapid population growth?} \hfill [Tag: Impacts of human population growth on environment | Easy]
\begin{enumerate}[label=(\alph*),topsep=2pt,itemsep=1pt,partopsep=0pt,parsep=0pt]
    \item {Lower generation of waste}
    \item {Improved forest coverage}
    \item {Reduced demand for land}
    \item {Increased pressure on resources}
\end{enumerate}
\vspace{2mm}

\noindent\textbf{Q4. How can population growth contribute to deforestation?} \hfill [Tag: Impacts of human population growth on environment | Easy]
\begin{enumerate}[label=(\alph*),topsep=2pt,itemsep=1pt,partopsep=0pt,parsep=0pt]
    \item {By limiting the use of timber}
    \item {By increasing demand for land}
    \item {By decreasing urban development}
    \item {By reducing the need for food}
\end{enumerate}
\vspace{2mm}

\noindent\textbf{Q5. Which measure helps prevent water-borne diseases?} \hfill [Tag: Human health and welfare | Easy]
\begin{enumerate}[label=(\alph*),topsep=2pt,itemsep=1pt,partopsep=0pt,parsep=0pt]
    \item {Drinking untreated water}
    \item {Using clean drinking water}
    \item {Storing waste near wells}
    \item {Allowing open defecation}
\end{enumerate}
\vspace{2mm}

\noindent\textbf{Q6. Poor air quality most directly affects which body system?} \hfill [Tag: Human health and welfare | Easy]
\begin{enumerate}[label=(\alph*),topsep=2pt,itemsep=1pt,partopsep=0pt,parsep=0pt]
    \item {Reproductive system}
    \item {Skeletal system}
    \item {Digestive system}
    \item {Respiratory system}
\end{enumerate}
\vspace{2mm}

\noindent\textbf{Q7. What is the central concern of environmental ethics?} \hfill [Tag: Environmental ethics | Easy]
\begin{enumerate}[label=(\alph*),topsep=2pt,itemsep=1pt,partopsep=0pt,parsep=0pt]
    \item {Methods of industrial accounting}
    \item {Rules for international trade}
    \item {Relations between humans and nature}
    \item {Techniques for urban construction}
\end{enumerate}
\vspace{2mm}

\noindent\textbf{Q8. What is a major purpose of environmental awareness campaigns?} \hfill [Tag: Public awareness | Easy]
\begin{enumerate}[label=(\alph*),topsep=2pt,itemsep=1pt,partopsep=0pt,parsep=0pt]
    \item {To discourage community participation}
    \item {To increase unnecessary resource consumption}
    \item {To encourage responsible environmental action}
    \item {To promote uncontrolled waste disposal}
\end{enumerate}
\vspace{2mm}

\noindent\textbf{Q9. How do sacred groves support environmental conservation?} \hfill [Tag: Cultural practices for environmental conservation | Easy]
\begin{enumerate}[label=(\alph*),topsep=2pt,itemsep=1pt,partopsep=0pt,parsep=0pt]
    \item {They promote large-scale hunting}
    \item {They protect patches of vegetation}
    \item {They replace forests with farms}
    \item {They provide sites for mining}
\end{enumerate}
\vspace{2mm}

\noindent\textbf{Q10. Which practice often helps tribal communities conserve natural resources?} \hfill [Tag: Tribal populations and environmental conservation | Easy]
\begin{enumerate}[label=(\alph*),topsep=2pt,itemsep=1pt,partopsep=0pt,parsep=0pt]
    \item {Unregulated industrial mining}
    \item {Traditional ecological knowledge}
    \item {Excessive use of chemicals}
    \item {Complete removal of forests}
\end{enumerate}
\vspace{2mm}

\noindent\textbf{Q11. Which Indian law was enacted specifically to protect wild animals, birds, and plants?} \hfill [Tag: Environmental laws in India | Easy]
\begin{enumerate}[label=(\alph*),topsep=2pt,itemsep=1pt,partopsep=0pt,parsep=0pt]
    \item {Forest Act, 1927}
    \item {Water Act, 1974}
    \item {Wildlife Protection Act, 1972}
    \item {Environment Protection Act, 1986}
\end{enumerate}
\vspace{2mm}

\noindent\textbf{Q12. The Water (Prevention and Control of Pollution) Act was enacted in which year?} \hfill [Tag: Environmental laws in India | Easy]
\begin{enumerate}[label=(\alph*),topsep=2pt,itemsep=1pt,partopsep=0pt,parsep=0pt]
    \item {1986}
    \item {1972}
    \item {1974}
    \item {1981}
\end{enumerate}
\vspace{2mm}

\noindent\textbf{Q13. Which law provides a broad framework for environmental protection in India?} \hfill [Tag: Environmental laws in India | Easy]
\begin{enumerate}[label=(\alph*),topsep=2pt,itemsep=1pt,partopsep=0pt,parsep=0pt]
    \item {Indian Forest Act, 1927}
    \item {Wildlife Protection Act, 1972}
    \item {Water Act, 1974}
    \item {Environment Protection Act, 1986}
\end{enumerate}
\vspace{2mm}

\noindent\textbf{Q14. Who led the Bishnoi sacrifice to protect Khejri trees at Khejarli?} \hfill [Tag: Environmental movements: Bishnois of Rajasthan | Easy]
\begin{enumerate}[label=(\alph*),topsep=2pt,itemsep=1pt,partopsep=0pt,parsep=0pt]
    \item {Medha Patkar}
    \item {Amrita Devi}
    \item {Sugathakumari}
    \item {Gaura Devi}
\end{enumerate}
\vspace{2mm}

\noindent\textbf{Q15. What action became the main symbol of the Chipko movement?} \hfill [Tag: Chipko movement | Easy]
\begin{enumerate}[label=(\alph*),topsep=2pt,itemsep=1pt,partopsep=0pt,parsep=0pt]
    \item {Cleaning rivers near cities}
    \item {Hugging trees to prevent felling}
    \item {Planting crops around dams}
    \item {Relocating wildlife from forests}
\end{enumerate}
\vspace{2mm}

\noindent\textbf{Q16. In which Indian state did the Appiko movement begin?} \hfill [Tag: Appiko movement | Easy]
\begin{enumerate}[label=(\alph*),topsep=2pt,itemsep=1pt,partopsep=0pt,parsep=0pt]
    \item {Karnataka}
    \item {Kerala}
    \item {Odisha}
    \item {Gujarat}
\end{enumerate}
\vspace{2mm}

\noindent\textbf{Q17. What was a major concern of the Jungle Bachao Andolan?} \hfill [Tag: Jungle Bachao Andolan | Easy]
\begin{enumerate}[label=(\alph*),topsep=2pt,itemsep=1pt,partopsep=0pt,parsep=0pt]
    \item {Building airports within major cities}
    \item {Replacing natural forests with teak plantations}
    \item {Increasing tourism in desert regions}
    \item {Expanding fishing in coastal waters}
\end{enumerate}
\vspace{2mm}

\noindent\textbf{Q18. The Silent Valley movement opposed a proposed project in which state?} \hfill [Tag: Silent Valley movement | Easy]
\begin{enumerate}[label=(\alph*),topsep=2pt,itemsep=1pt,partopsep=0pt,parsep=0pt]
    \item {Bihar}
    \item {Kerala}
    \item {Rajasthan}
    \item {Punjab}
\end{enumerate}
\vspace{2mm}

\noindent\textbf{Q19. Which issue was central to the Narmada Bachao Andolan?} \hfill [Tag: Narmada Bachao Andolan | Easy]
\begin{enumerate}[label=(\alph*),topsep=2pt,itemsep=1pt,partopsep=0pt,parsep=0pt]
    \item {Protection of monuments from air pollution}
    \item {Control of traffic in crowded cities}
    \item {Impacts of large dams on people and nature}
    \item {Expansion of industries in dry regions}
\end{enumerate}
\vspace{2mm}

\noindent\textbf{Q20. The Tehri Dam was constructed on which river?} \hfill [Tag: Tehri Dam conflict | Easy]
\begin{enumerate}[label=(\alph*),topsep=2pt,itemsep=1pt,partopsep=0pt,parsep=0pt]
    \item {Mahanadi}
    \item {Bhagirathi}
    \item {Godavari}
    \item {Yamuna}
\end{enumerate}
\vspace{2mm}

\noindent\textbf{Q21. A town's population doubles in 35 years while maintaining an approximately constant exponential growth rate. What is its approximate annual growth rate?} \hfill [Tag: Human population growth | Medium]
\begin{enumerate}[label=(\alph*),topsep=2pt,itemsep=1pt,partopsep=0pt,parsep=0pt]
    \item {1.0\% per year}
    \item {5.0\% per year}
    \item {2.0\% per year}
    \item {3.5\% per year}
\end{enumerate}
\vspace{2mm}

\noindent\textbf{Q22. A country has a population pyramid with a very broad base and a narrow top. Which policy need is most likely to become urgent over the next decade?} \hfill [Tag: Human population growth | Medium]
\begin{enumerate}[label=(\alph*),topsep=2pt,itemsep=1pt,partopsep=0pt,parsep=0pt]
    \item {Closing maternity and childcare facilities}
    \item {Expanding schools and employment opportunities}
    \item {Reducing pensions for an ageing population}
    \item {Increasing immigration of young workers}
\end{enumerate}
\vspace{2mm}

\noindent\textbf{Q23. A rapidly growing city fills nearby wetlands to construct housing. During the following monsoon, urban flooding becomes more frequent. What best explains this change?} \hfill [Tag: Impacts of human population growth on environment | Medium]
\begin{enumerate}[label=(\alph*),topsep=2pt,itemsep=1pt,partopsep=0pt,parsep=0pt]
    \item {Wetland loss increased groundwater recharge}
    \item {Housing reduced the volume of surface runoff}
    \item {Housing increased evaporation from the soil}
    \item {Wetland loss reduced natural water storage}
\end{enumerate}
\vspace{2mm}

\noindent\textbf{Q24. Population and water demand are rising in a semi-arid district where groundwater levels are declining. Which response addresses both consumption and recharge?} \hfill [Tag: Impacts of human population growth on environment | Medium]
\begin{enumerate}[label=(\alph*),topsep=2pt,itemsep=1pt,partopsep=0pt,parsep=0pt]
    \item {Deepening every private borewell}
    \item {Importing bottled water for households}
    \item {Installing meters and restoring recharge ponds}
    \item {Lining all open land with concrete}
\end{enumerate}
\vspace{2mm}

\noindent\textbf{Q25. After heavy rainfall, a settlement reports many cases of diarrhoea. Tests show that sewage has entered drinking-water wells. Which intervention should receive immediate priority?} \hfill [Tag: Human health and welfare | Medium]
\begin{enumerate}[label=(\alph*),topsep=2pt,itemsep=1pt,partopsep=0pt,parsep=0pt]
    \item {Increasing electricity supply to homes}
    \item {Disinfecting water and repairing sanitation}
    \item {Expanding roads around the settlement}
    \item {Planting shade trees near the wells}
\end{enumerate}
\vspace{2mm}

\noindent\textbf{Q26. During a severe urban heat wave, which measure would most directly protect vulnerable residents while also improving long-term welfare?} \hfill [Tag: Human health and welfare | Medium]
\begin{enumerate}[label=(\alph*),topsep=2pt,itemsep=1pt,partopsep=0pt,parsep=0pt]
    \item {Widening highways and increasing parking spaces}
    \item {Opening cooling centres and expanding tree cover}
    \item {Replacing reflective roofs with dark roofs}
    \item {Closing public parks during afternoon hours}
\end{enumerate}
\vspace{2mm}

\noindent\textbf{Q27. A mining project will provide income for 15 years but may leave toxic waste requiring management for centuries. Which ethical principle most strongly supports considering those future costs?} \hfill [Tag: Environmental ethics | Medium]
\begin{enumerate}[label=(\alph*),topsep=2pt,itemsep=1pt,partopsep=0pt,parsep=0pt]
    \item {Administrative efficiency}
    \item {Consumer sovereignty}
    \item {Market competition}
    \item {Intergenerational equity}
\end{enumerate}
\vspace{2mm}

\noindent\textbf{Q28. A pesticide may seriously harm pollinators, but available scientific evidence is incomplete. Which decision best reflects the precautionary principle?} \hfill [Tag: Environmental ethics | Medium]
\begin{enumerate}[label=(\alph*),topsep=2pt,itemsep=1pt,partopsep=0pt,parsep=0pt]
    \item {Restrict use while conducting further research}
    \item {Ignore ecological effects if crop yields increase}
    \item {Delay monitoring until pollinator numbers decline}
    \item {Permit unrestricted use until damage is proven}
\end{enumerate}
\vspace{2mm}

\noindent\textbf{Q29. A municipality wants an air-pollution awareness campaign to produce measurable behavioural change. Which approach is most suitable?} \hfill [Tag: Public awareness | Medium]
\begin{enumerate}[label=(\alph*),topsep=2pt,itemsep=1pt,partopsep=0pt,parsep=0pt]
    \item {Publishing one annual report for technical experts}
    \item {Using general slogans without identifying sources}
    \item {Sharing local alerts with transport alternatives}
    \item {Displaying pollution data without practical advice}
\end{enumerate}
\vspace{2mm}

\noindent\textbf{Q30. A village protects a forest patch because it is considered sacred and prohibits tree felling there. What is the main conservation value of this practice?} \hfill [Tag: Cultural practices for environmental conservation | Medium]
\begin{enumerate}[label=(\alph*),topsep=2pt,itemsep=1pt,partopsep=0pt,parsep=0pt]
    \item {It guarantees unrestricted extraction of forest produce}
    \item {It eliminates the need to conserve nearby ecosystems}
    \item {It preserves habitat through community restrictions}
    \item {It converts native forest into commercial farmland}
\end{enumerate}
\vspace{2mm}

\noindent\textbf{Q31. A forest department plans habitat restoration in an area traditionally managed by a tribal community. Which strategy is most likely to produce durable conservation outcomes?} \hfill [Tag: Tribal populations and environmental conservation | Medium]
\begin{enumerate}[label=(\alph*),topsep=2pt,itemsep=1pt,partopsep=0pt,parsep=0pt]
    \item {Combining local knowledge with ecological research}
    \item {Excluding residents from all planning meetings}
    \item {Replacing local species with commercial trees}
    \item {Relocating residents before assessing their practices}
\end{enumerate}
\vspace{2mm}

\noindent\textbf{Q32. A state agency proposes using reserved forest land for a highway. Under the Forest (Conservation) Act, 1980, what is a key legal requirement?} \hfill [Tag: Environmental laws in India | Medium]
\begin{enumerate}[label=(\alph*),topsep=2pt,itemsep=1pt,partopsep=0pt,parsep=0pt]
    \item {Approval only from the nearest municipality}
    \item {Prior approval from the Central Government}
    \item {Automatic clearance after a public notice}
    \item {Consent only from the highway contractor}
\end{enumerate}
\vspace{2mm}

\noindent\textbf{Q33. An industrial accident releases a hazardous chemical not adequately addressed by a narrowly focused pollution rule. Which Indian law gives the Central Government broad powers to coordinate environmental protection?} \hfill [Tag: Environmental laws in India | Medium]
\begin{enumerate}[label=(\alph*),topsep=2pt,itemsep=1pt,partopsep=0pt,parsep=0pt]
    \item {Indian Forest Act, 1927}
    \item {Biological Diversity Act, 2002}
    \item {Environment (Protection) Act, 1986}
    \item {Wildlife (Protection) Act, 1972}
\end{enumerate}
\vspace{2mm}

\noindent\textbf{Q34. A community in arid Rajasthan protects Khejri trees and wild animals despite economic pressure to exploit them. Which historical conservation tradition does this most closely reflect?} \hfill [Tag: Environmental movements: Bishnois of Rajasthan | Medium]
\begin{enumerate}[label=(\alph*),topsep=2pt,itemsep=1pt,partopsep=0pt,parsep=0pt]
    \item {The Appiko tradition}
    \item {The Silent Valley campaign}
    \item {The Bishnoi tradition}
    \item {The Narmada campaign}
\end{enumerate}
\vspace{2mm}

\noindent\textbf{Q35. Villagers in a Himalayan region embrace trees to prevent commercial logging because the forest supports fuel, fodder, soil stability, and water security. Which movement does this action represent?} \hfill [Tag: Chipko movement | Medium]
\begin{enumerate}[label=(\alph*),topsep=2pt,itemsep=1pt,partopsep=0pt,parsep=0pt]
    \item {Jungle Bachao Andolan}
    \item {Appiko movement}
    \item {Chipko movement}
    \item {Silent Valley movement}
\end{enumerate}
\vspace{2mm}

\noindent\textbf{Q36. A Western Ghats campaign in Karnataka promotes conserving existing forests, regenerating degraded land, and using forest resources carefully. Which movement matches these goals?} \hfill [Tag: Appiko movement | Medium]
\begin{enumerate}[label=(\alph*),topsep=2pt,itemsep=1pt,partopsep=0pt,parsep=0pt]
    \item {Chipko movement}
    \item {Appiko movement}
    \item {Tehri Dam conflict}
    \item {Narmada Bachao Andolan}
\end{enumerate}
\vspace{2mm}

\noindent\textbf{Q37. A forest-dependent community opposes replacing biologically diverse sal forests with commercially valuable teak plantations. Which movement is most directly associated with such a protest?} \hfill [Tag: Jungle Bachao Andolan | Medium]
\begin{enumerate}[label=(\alph*),topsep=2pt,itemsep=1pt,partopsep=0pt,parsep=0pt]
    \item {Silent Valley movement}
    \item {Narmada Bachao Andolan}
    \item {Jungle Bachao Andolan}
    \item {Chipko movement}
\end{enumerate}
\vspace{2mm}

\noindent\textbf{Q38. A proposed hydroelectric project threatens an undisturbed tropical evergreen forest in Kerala that supports species such as the lion-tailed macaque. Which campaign provides the closest historical parallel?} \hfill [Tag: Silent Valley movement | Medium]
\begin{enumerate}[label=(\alph*),topsep=2pt,itemsep=1pt,partopsep=0pt,parsep=0pt]
    \item {Appiko movement}
    \item {Tehri Dam conflict}
    \item {Bishnoi movement}
    \item {Silent Valley movement}
\end{enumerate}
\vspace{2mm}

\noindent\textbf{Q39. Residents affected by a large river-valley project demand proper rehabilitation and question the social and ecological costs of multiple large dams. Which movement best represents these concerns?} \hfill [Tag: Narmada Bachao Andolan | Medium]
\begin{enumerate}[label=(\alph*),topsep=2pt,itemsep=1pt,partopsep=0pt,parsep=0pt]
    \item {Chipko movement}
    \item {Appiko movement}
    \item {Narmada Bachao Andolan}
    \item {Jungle Bachao Andolan}
\end{enumerate}
\vspace{2mm}

\noindent\textbf{Q40. A debate concerns a major dam on the Bhagirathi River, the displacement of local communities, and safety risks in an earthquake-prone Himalayan region. Which conflict is being described?} \hfill [Tag: Tehri Dam conflict | Medium]
\begin{enumerate}[label=(\alph*),topsep=2pt,itemsep=1pt,partopsep=0pt,parsep=0pt]
    \item {Jungle Bachao conflict}
    \item {Tehri Dam conflict}
    \item {Appiko conflict}
    \item {Silent Valley conflict}
\end{enumerate}
\vspace{2mm}

\noindent\textbf{Q41. A country reaches replacement-level fertility, but its population continues to grow for several decades. Which demographic mechanism best explains this outcome?} \hfill [Tag: Human population growth | Hard]
\begin{enumerate}[label=(\alph*),topsep=2pt,itemsep=1pt,partopsep=0pt,parsep=0pt]
    \item {Life expectancy falls while the crude death rate remains constant}
    \item {The sex ratio shifts toward older women outside reproductive age}
    \item {Net migration becomes zero after fertility reaches replacement level}
    \item {A large proportion of the population is entering reproductive age}
\end{enumerate}
\vspace{2mm}

\noindent\textbf{Q42. A population of 2 million has annual birth, death, immigration, and emigration rates of 24, 8, 3, and 5 per 1,000 people, respectively. Assuming these rates remain constant, what is its approximate net numerical change during one year?} \hfill [Tag: Human population growth | Hard]
\begin{enumerate}[label=(\alph*),topsep=2pt,itemsep=1pt,partopsep=0pt,parsep=0pt]
    \item {An increase of 40,000 people}
    \item {An increase of 28,000 people}
    \item {An increase of 20,000 people}
    \item {An increase of 32,000 people}
\end{enumerate}
\vspace{2mm}

\noindent\textbf{Q43. Using the relation I = P $\times$ A $\times$ T , population rises by 20\%, affluence per person rises by 25\%, and environmental impact per unit consumption falls by 20\%. What is the resulting change in total environmental impact?} \hfill [Tag: Impacts of human population growth on environment | Hard]
\begin{enumerate}[label=(\alph*),topsep=2pt,itemsep=1pt,partopsep=0pt,parsep=0pt]
    \item {A 4\% decrease}
    \item {No net change}
    \item {A 20\% increase}
    \item {A 25\% increase}
\end{enumerate}
\vspace{2mm}

\noindent\textbf{Q44. Two cities have equal populations. City X has high-density housing and mass transit, while City Y has dispersed housing, private vehicles, and fragmented infrastructure. Which conclusion is most defensible?} \hfill [Tag: Impacts of human population growth on environment | Hard]
\begin{enumerate}[label=(\alph*),topsep=2pt,itemsep=1pt,partopsep=0pt,parsep=0pt]
    \item {City Y must have lower impact because its pollution sources are geographically dispersed}
    \item {City Y may have greater per-capita impact because settlement form affects consumption}
    \item {Equal populations necessarily produce equal ecological footprints}
    \item {City X must have greater impact because density increases local exposure}
\end{enumerate}
\vspace{2mm}

\noindent\textbf{Q45. An intervention reduces deaths from indoor air pollution but leaves many survivors with chronic respiratory disability. Which metric is best suited to evaluating both effects together?} \hfill [Tag: Human health and welfare | Hard]
\begin{enumerate}[label=(\alph*),topsep=2pt,itemsep=1pt,partopsep=0pt,parsep=0pt]
    \item {Disability-adjusted life years}
    \item {Crude birth rate}
    \item {Infant mortality rate}
    \item {Point prevalence alone}
\end{enumerate}
\vspace{2mm}

\noindent\textbf{Q46. After flooding, a settlement experiences diarrhoeal disease, mosquito proliferation, and loss of access to health services. Which response most directly reflects an integrated environmental-health strategy?} \hfill [Tag: Human health and welfare | Hard]
\begin{enumerate}[label=(\alph*),topsep=2pt,itemsep=1pt,partopsep=0pt,parsep=0pt]
    \item {Relocate only confirmed patients while leaving water systems unchanged}
    \item {Distribute antibiotics to all residents regardless of diagnosis}
    \item {Restore safe water, sanitation, vector control, and primary care access}
    \item {Begin mosquito control while postponing sanitation restoration}
\end{enumerate}
\vspace{2mm}

\noindent\textbf{Q47. A forest project is rejected because an endemic species is considered valuable even if it provides no known benefit to humans. Which ethical position most directly supports the decision?} \hfill [Tag: Environmental ethics | Hard]
\begin{enumerate}[label=(\alph*),topsep=2pt,itemsep=1pt,partopsep=0pt,parsep=0pt]
    \item {Strict anthropocentrism}
    \item {Market utilitarianism}
    \item {Instrumental conservation}
    \item {Biocentric intrinsic value}
\end{enumerate}
\vspace{2mm}

\noindent\textbf{Q48. A government permits depletion of a unique wetland because the revenue will finance roads and schools of equivalent monetary value. Which principle most strongly challenges the assumption that this substitution is ethically sufficient?} \hfill [Tag: Environmental ethics | Hard]
\begin{enumerate}[label=(\alph*),topsep=2pt,itemsep=1pt,partopsep=0pt,parsep=0pt]
    \item {Weak sustainability}
    \item {Consumer sovereignty}
    \item {Strong sustainability}
    \item {Hedonic pricing}
\end{enumerate}
\vspace{2mm}

\noindent\textbf{Q49. An awareness campaign raises knowledge about household waste segregation, but observed segregation rates remain unchanged. Which redesign is most likely to close this knowledge-action gap?} \hfill [Tag: Public awareness | Hard]
\begin{enumerate}[label=(\alph*),topsep=2pt,itemsep=1pt,partopsep=0pt,parsep=0pt]
    \item {Combine clear prompts with separate bins and reliable collection}
    \item {Measure campaign reach using advertisement impressions alone}
    \item {Provide more technical facts without changing collection arrangements}
    \item {Replace local messages with general global pollution statistics}
\end{enumerate}
\vspace{2mm}

\noindent\textbf{Q50. A sacred grove protects old-growth vegetation but excludes customary collection by some marginalized residents. Which assessment best integrates conservation and equity?} \hfill [Tag: Cultural practices for environmental conservation | Hard]
\begin{enumerate}[label=(\alph*),topsep=2pt,itemsep=1pt,partopsep=0pt,parsep=0pt]
    \item {Preserve ecological protection while revising inequitable governance rules}
    \item {Replace community governance with unrestricted commercial resource access}
    \item {Abolish the grove because any exclusion invalidates conservation}
    \item {Retain all customary rules because traditional practices are inherently equitable}
\end{enumerate}
\vspace{2mm}

\noindent\textbf{Q51. A tribal community manages a forest through seasonal harvest restrictions, habitat taboos, and collective monitoring. A conservation plan instead removes the community and centralizes enforcement. What key risk does the plan create?} \hfill [Tag: Tribal populations and environmental conservation | Hard]
\begin{enumerate}[label=(\alph*),topsep=2pt,itemsep=1pt,partopsep=0pt,parsep=0pt]
    \item {It guarantees that all commercial extraction will immediately cease}
    \item {It may remove local institutions and knowledge that regulate resource use}
    \item {It converts customary knowledge into stronger statutory protection automatically}
    \item {It ensures that ecological monitoring becomes more spatially detailed}
\end{enumerate}
\vspace{2mm}

\noindent\textbf{Q52. A company begins constructing an industrial unit without a required prior environmental clearance and argues that compensatory measures can be imposed after construction. Which legal principle most directly undermines this argument?} \hfill [Tag: Environmental laws in India | Hard]
\begin{enumerate}[label=(\alph*),topsep=2pt,itemsep=1pt,partopsep=0pt,parsep=0pt]
    \item {State land ownership eliminates the need for environmental appraisal}
    \item {Prior appraisal is preventive and cannot be reduced to routine post-facto approval}
    \item {Compensation automatically substitutes for every procedural requirement}
    \item {Environmental clearance is required only after commercial production begins}
\end{enumerate}
\vspace{2mm}

\noindent\textbf{Q53. A proposal requires diversion of reserved forest land for a non-forest purpose and may also affect habitat of a protected wild animal. Which combination of statutes is most directly relevant?} \hfill [Tag: Environmental laws in India | Hard]
\begin{enumerate}[label=(\alph*),topsep=2pt,itemsep=1pt,partopsep=0pt,parsep=0pt]
    \item {Forest (Conservation) Act, 1980 and Wildlife (Protection) Act, 1972}
    \item {National Green Tribunal Act, 2010 and Indian Forest Act, 1927 only}
    \item {Water Act, 1974 and Air Act, 1981}
    \item {Biological Diversity Act, 2002 and Public Liability Insurance Act, 1991}
\end{enumerate}
\vspace{2mm}

\noindent\textbf{Q54. A foreign company seeks access to an Indian biological resource for commercial research, while local knowledge holders claim a share of resulting benefits. Which statute most directly provides the access-and-benefit-sharing framework?} \hfill [Tag: Environmental laws in India | Hard]
\begin{enumerate}[label=(\alph*),topsep=2pt,itemsep=1pt,partopsep=0pt,parsep=0pt]
    \item {Biological Diversity Act, 2002}
    \item {Forest (Conservation) Act, 1980}
    \item {Water (Prevention and Control of Pollution) Act, 1974}
    \item {Environment (Protection) Act, 1986}
\end{enumerate}
\vspace{2mm}

\noindent\textbf{Q55. Which interpretation best explains why the 1730 Khejarli sacrifice is significant in environmental history?} \hfill [Tag: Environmental movements: Bishnois of Rajasthan | Hard]
\begin{enumerate}[label=(\alph*),topsep=2pt,itemsep=1pt,partopsep=0pt,parsep=0pt]
    \item {It linked defense of khejri trees to a community's religious and livelihood ethics}
    \item {It demanded replacement of commercial teak plantations with natural sal forests}
    \item {It opposed a hydroelectric project primarily because of seismic hazards}
    \item {It created India's first statutory national park through judicial intervention}
\end{enumerate}
\vspace{2mm}

\noindent\textbf{Q56. Which comparison correctly distinguishes the Reni episode of the Chipko movement from a narrow timber-management dispute?} \hfill [Tag: Environmental movements: Chipko movement | Hard]
\begin{enumerate}[label=(\alph*),topsep=2pt,itemsep=1pt,partopsep=0pt,parsep=0pt]
    \item {It sought to replace mixed Himalayan forests with fast-growing monocultures}
    \item {It focused solely on royalties paid to contractors for exported timber}
    \item {It opposed wildlife protection because it restricted industrial logging}
    \item {It connected women's forest dependence with resistance to commercial felling}
\end{enumerate}
\vspace{2mm}

\noindent\textbf{Q57. Which strategy most accurately reflects the broader program associated with the Appiko movement in Karnataka?} \hfill [Tag: Environmental movements: Appiko movement | Hard]
\begin{enumerate}[label=(\alph*),topsep=2pt,itemsep=1pt,partopsep=0pt,parsep=0pt]
    \item {Protect plantations, expand clear-felling, and centralize forest-product markets}
    \item {Protect existing forests, regenerate degraded areas, and use resources rationally}
    \item {Nationalize sacred groves, prohibit restoration, and maximize timber auctions}
    \item {Stop irrigation projects, remove local access, and promote mining employment}
\end{enumerate}
\vspace{2mm}

\noindent\textbf{Q58. The Jungle Bachao Andolan in the Singhbhum region is best understood as resistance to which forestry transition?} \hfill [Tag: Environmental movements: Jungle Bachao Andolan | Hard]
\begin{enumerate}[label=(\alph*),topsep=2pt,itemsep=1pt,partopsep=0pt,parsep=0pt]
    \item {Replacement of natural sal forests with commercial teak plantations}
    \item {Conversion of alpine meadows into state-owned apple orchards}
    \item {Replacement of tea estates with community-managed mangrove forests}
    \item {Conversion of desert commons into khejri-based agroforestry systems}
\end{enumerate}
\vspace{2mm}

\noindent\textbf{Q59. Why did the proposed Silent Valley hydroelectric project become a national conservation issue despite being presented as a renewable-energy project?} \hfill [Tag: Environmental movements: Silent Valley movement | Hard]
\begin{enumerate}[label=(\alph*),topsep=2pt,itemsep=1pt,partopsep=0pt,parsep=0pt]
    \item {Its transmission lines crossed no forest but displaced an entire coastal fishery}
    \item {Its canals would have replaced Rajasthan's desert forests with teak plantations}
    \item {Its turbines required continuous combustion of imported high-sulfur coal}
    \item {Its reservoir threatened an exceptional tropical evergreen ecosystem and endemic fauna}
\end{enumerate}
\vspace{2mm}

\noindent\textbf{Q60. Which policy conflict lies most centrally at the heart of the Narmada Bachao Andolan's challenge to large dams such as Sardar Sarovar?} \hfill [Tag: Environmental movements: Narmada Bachao Andolan | Hard]
\begin{enumerate}[label=(\alph*),topsep=2pt,itemsep=1pt,partopsep=0pt,parsep=0pt]
    \item {Distribution of development benefits versus displacement and ecological costs}
    \item {Commercial tourism growth versus access to Himalayan pilgrimage routes}
    \item {Urban vehicle ownership versus enforcement of national fuel standards}
    \item {Expansion of marine fisheries versus protection of offshore coral reefs}
\end{enumerate}
\vspace{2mm}


\section{Section B: Comprehensive Subjective Problem with Detailed Analytical Rubric}

\begin{subjectiveq}{6}{1}
\textbf{Socio-Ecological Conflicts, Constitutional Jurisprudence and Movement Dynamics:}\\A multi-purpose hydro-electric river valley project threatens to submerge $35{,}000\text{ hectares}$ of pristine riparian sal forest and displace $40{,}000$ indigenous tribal families residing across $180$ villages in central India.\\\\(a) Examine the ideological conflict between Anthropocentrism (instrumental human-centric resource extraction) and Ecocentrism/Biocentrism (intrinsic moral standing of non-human entities and ecosystems).\\(b) Detail the Indian constitutional mandate protecting the environment, specifically analyzing the intersection of Article 21 (Fundamental Right to Life including a clean, pollution-free environment), Article 48A (Directive Principles of State Policy), and Article 51A(g) (Fundamental Duties of citizens).\\(c) Compare the historical paradigms, leadership structures, mobilization tactics, and outcomes of: (i) The Chipko Movement (1973, Uttarakhand), (ii) The Silent Valley Movement (1978, Kerala), and (iii) The Narmada Bachao Andolan (1985--present).\\(d) Outline the statutory provisions of the Scheduled Tribes and Other Traditional Forest Dwellers (Recognition of Forest Rights) Act, 2006 (FRA) and the Environmental Impact Assessment (EIA) Notification regarding mandatory public consultation.
\end{subjectiveq}

\begin{solbox}
\textbf{Part (a): Anthropocentrism versus Ecocentrism:}\\\begin{itemize}[leftmargin=1.5em]  \item \textbf{Anthropocentrism:} Evaluates ecosystems strictly as natural capital and instrumental resource pools designed to maximize human economic utility, industrial GDP growth, and power generation.  \item \textbf{Ecocentrism and Deep Ecology (Arne Naess):} Recognizes that biotic communities, rivers, and wildlife possess inherent, non-negotiable moral standing independent of human utility. Ecosystem integrity must be preserved for intergenerational and inter-species justice.\end{itemize}\textbf{Part (b): Indian Constitutional Environmental Jurisprudence:}\\The Supreme Court of India, through landmark public interest litigation (PIL) rulings (\textit{Subhash Kumar v. State of Bihar, 1991}; \textit{Vellore Citizens Welfare Forum, 1996}; \textit{M.C. Mehta series}), expanded environmental rights:\begin{itemize}[leftmargin=1.5em]  \item \textbf{Article 21 (Fundamental Right to Life):} Interpreted expansively to guarantee the right to wholesome, unpolluted air, clean drinking water, and ecological dignity.  \item \textbf{Article 48A (Directive Principles of State Policy, 42nd Amendment 1976):} Mandates that the State shall endeavor to protect and improve the environment and safeguard forests and wildlife.  \item \textbf{Article 51A(g) (Fundamental Duties):} Imposes a positive constitutional duty upon every Indian citizen to protect and improve the natural environment including forests, lakes, rivers, and wildlife, and to have compassion for living creatures.\end{itemize}\textbf{Part (c): Comparative Movement Analysis:}\\\begin{center}\begin{tabularx}{\textwidth}{p{2.8cm} X X X}\toprule\textbf{Dimension} & \textbf{Chipko Movement (1973)} & \textbf{Silent Valley (1978)} & \textbf{Narmada Bachao (1985)} \\\midrule \textbf{Geographic Core} & Garhwal Himalayas (Reni, Chamoli, Uttarakhand) & Palakkad district, Western Ghats, Kerala & Narmada river valley (MP, Maharashtra, Gujarat) \\\textbf{Key Leadership} & Gaura Devi, Sunderlal Bahuguna, Chandi Prasad Bhatt & Kerala Sastra Sahitya Parishad (KSSP), Salim Ali & Medha Patkar, Baba Amte, indigenous Adivasi elders \\\textbf{Tactical Focus} & Hugging trees to halt commercial logging auctions & Mass public science education, civil society \& scientist petitions & Non-violent satyagraha, jal-samarpan, hunger strikes, legal PILs \\\textbf{Concrete Outcome} & 15-year commercial green tree felling ban in Himalayan forests (1980) & Hydro-dam cancelled; declared Silent Valley National Park (1984) & Landmark rehabilitation judgments; international World Bank loan withdrawal \\\bottomrule\end{tabularx}\end{center}\textbf{Part (d): Statutory Safeguards:}\\Under the \textbf{Forest Rights Act (FRA) 2006}, the Gram Sabha holds absolute legal authority to grant or withhold free, prior, and informed consent (FPIC) before any community forest resource is diverted for non-forest infrastructure, as affirmed in the historic Niyamgiri judgment (2013). Under the \textbf{EIA Notification 2006}, Category A mega-projects require mandatory public hearings conducted by the State Pollution Control Board.
\end{solbox}

\begin{examinertip}
\textbf{Marking Scheme Breakdown (Total: 10 Marks):}\\
\textbullet\ \textbf{3 Marks:} Clear formulation of governing environmental identities, equations, and thermodynamic/constitutional foundations.\\
\textbullet\ \textbf{4 Marks:} Rigorous algebraic/numerical substitution, dimensional consistency, and unit handling.\\
\textbullet\ \textbf{3 Marks:} Critical evaluation of environmental implications, regulatory criteria (CPCB/UNDRR/IUCN), and policy mitigations.
\end{examinertip}
